\documentclass[10pt]{beamer}

\usetheme{Madrid}
\usecolortheme{whale}

\usepackage[utf8]{inputenc}
\usepackage[T1]{fontenc}
\usepackage[italian]{babel}
\usepackage{tikz}
\usetikzlibrary{arrows.meta, positioning, shapes.geometric, calc, fit}
\usepackage{listings}
\usepackage{xcolor}
\usepackage{booktabs}

% --- Configurazione Colori e Link ---
\definecolor{primary}{RGB}{38, 66, 139}
\definecolor{codebg}{RGB}{245, 245, 245}
\definecolor{keywordcolor}{RGB}{0, 0, 255}
\definecolor{stringcolor}{RGB}{163, 21, 21}
\definecolor{commentcolor}{RGB}{0, 128, 0}

\lstdefinelanguage{javascript}{
    keywords={async, await, break, case, catch, class, const, continue, debugger, default, delete, do, else, export, extends, false, finally, for, function, if, import, in, instanceof, new, null, return, super, switch, this, throw, true, try, typeof, var, void, while, with, yield, let},
    morekeywords={console, log, Error, fetch, JSON, parse, stringify},
    sensitive=true,
    morecomment=[l]{//},
    morecomment=[s]{/*}{*/},
    morestring=[b]',
    morestring=[b]",
    morestring=[b]`,
    keywordstyle=\color{primary}\bfseries,
    commentstyle=\color{commentcolor}\itshape,
    stringstyle=\color{stringcolor},
    basicstyle=\ttfamily\small
}

\lstset{
    backgroundcolor=\color{codebg},
    basicstyle=\ttfamily\small,
    breaklines=true,
    frame=single,
    showstringspaces=false,
    keywordstyle=\color{blue},
    commentstyle=\color{green!50!black},
    stringstyle=\color{red!70!black}
}

\title[Lezione 5]{Programmazione e Tecnologie Web}
\subtitle{I Framework Frontend Moderni (CSR vs SSR)}
\author{Giuseppe Capasso}
\date{2025/2026}

\begin{document}

\begin{frame}
    \titlepage
\end{frame}

\section{Evoluzione del Frontend}
\begin{frame}{Dal DOM Imperativo a React}
    Perché abbiamo smesso di manipolare il DOM direttamente?
    \begin{itemize}
        \item \textbf{Stato Frammentato}: Difficoltà nel sincronizzare dati e UI.
        \item \textbf{Performance}: Le operazioni sul DOM reale sono lente.
        \item \textbf{Manutenibilità}: Codice difficile da leggere e testare ("Spaghetti Code").
    \end{itemize}
    \vspace{0.3cm}
    \textbf{La Soluzione Dichiarativa}:
    \textit{"Descrivi COSA vuoi vedere, non COME farlo."}
\end{frame}

\begin{frame}{Il Virtual DOM}
    \centering
    \begin{tikzpicture}[
        node distance=1cm,
        box/.style={rectangle, draw, fill=blue!10, minimum width=2.5cm, minimum height=0.8cm, align=center, font=\small},
        arrow/.style={-latex, thick, primary}
    ]
        \node (state) [box, fill=yellow!10] {Stato\\(Dati)};
        \node (vdom) [box, right=1.5cm of state] {Virtual DOM\\(Memoria)};
        \node (realdom) [box, right=1.5cm of vdom, fill=green!10] {DOM Reale\\(Browser)};

        \draw [arrow] (state) -- node[above, scale=0.7] {Render} (vdom);
        \draw [arrow] (vdom) -- node[above, scale=0.7] {Diffing} (realdom);
    \end{tikzpicture}
    \begin{itemize}
        \item React calcola solo le differenze minime.
        \item Aggiornamenti chirurgici e ultra-veloci.
    \end{itemize}
\end{frame}

\section{Next.js Basics}
\begin{frame}[fragile]{Setup e Struttura}
    Next.js è il framework standard per applicazioni React professionali.
    \begin{itemize}
        \item \textbf{Installazione}: \texttt{npx create-next-app@latest}
        \item \textbf{App Router}: La cartella \texttt{app/} gestisce il routing.
    \end{itemize}
    \begin{block}{File Fondamentali}
        \begin{itemize}
            \item \texttt{page.js}: Il contenuto della pagina.
            \item \texttt{layout.js}: Elementi comuni (navbar, footer).
            \item \texttt{globals.css}: Stili globali e Tailwind.
        \end{itemize}
    \end{block}
\end{frame}

\begin{frame}[fragile]{Sintassi JSX}
    JSX permette di scrivere strutture simili ad HTML in JavaScript.
    \begin{lstlisting}[language=javascript]
export default function Home() {
  const user = "Giuseppe";
  return (
    <div className="p-4">
      <h1 className="text-xl">Ciao, {user}!</h1>
      {user === "Giuseppe" && <p>Docente</p>}
    </div>
  );
}
\end{lstlisting}
\end{frame}

\section{State Management}
\begin{frame}[fragile]{Gestione dello Stato: useState}
    Il componente reagisce ai cambiamenti dei dati.
    \begin{lstlisting}[language=javascript]
"use client";
import { useState } from "react";

const [search, setSearch] = useState("");

return (
  <input
    value={search}
    onChange={(e) => setSearch(e.target.value)}
  />
);
\end{lstlisting}
\end{frame}

\begin{frame}[fragile]{Stato Derivato e Filtraggio}
    Non creare stati inutili! Calcola i dati durante il rendering.
    \begin{lstlisting}[language=javascript]
// tasks e search sono variabili di stato
const filteredTasks = tasks.filter(t =>
  t.text.includes(search)
);

return (
  <ul>
    {filteredTasks.map(t => <li key={t.id}>{t.text}</li>)}
  </ul>
);
\end{lstlisting}
\end{frame}

\section{Sicurezza e Comunicazione}
\begin{frame}{Same-Origin Policy (SOP)}
    Il pilastro della sicurezza web nel browser.
    \begin{itemize}
        \item \textbf{Definizione}: Uno script può accedere a risorse solo se hanno la stessa \textit{Origine}.
        \item \textbf{Cosa definisce l'Origine?}
        \begin{enumerate}
            \item Protocollo (http vs https)
            \item Dominio (localhost vs google.com)
            \item Porta (3000 vs 5000)
        \end{enumerate}
    \end{itemize}
    \vspace{0.3cm}
    \begin{exampleblock}{Esempio di violazione}
        Frontend su \texttt{:3000} prova a chiamare API su \texttt{:5000}.
    \end{exampleblock}
\end{frame}

\begin{frame}{CORS: Cross-Origin Resource Sharing}
    Come "bypassare" la SOP in modo sicuro.
    \begin{itemize}
        \item Il server deve inviare l'header: \\ \texttt{Access-Control-Allow-Origin: *}
        \item \textbf{Preflight Request}: Per metodi "pericolosi" (POST, DELETE, ecc.), il browser invia prima una richiesta \textbf{OPTIONS}.
    \end{itemize}
    \centering
    \begin{tikzpicture}[
        node distance=3cm, 
        font=\tiny,
        box/.style={rectangle, draw, rounded corners, minimum width=2cm, minimum height=0.8cm, align=center, thick},
        arrow/.style={-latex, thick}
    ]
        \node (br) [box, fill=blue!10] {Browser\\(Client)};
        \node (srv) [box, right=of br, fill=green!10] {Server API\\(Backend)};

        % 1. Preflight
        \draw[arrow, primary] ([yshift=3mm]br.east) -- node[above, black] {1. OPTIONS (Preflight)} ([yshift=3mm]srv.west);
        \draw[arrow, gray, dashed] ([yshift=1mm]srv.west) -- node[below, black] {2. OK (CORS Headers)} ([yshift=1mm]br.east);

        % 2. Actual Request
        \draw[arrow, primary] ([yshift=-2mm]br.east) -- node[above, black] {3. Richiesta Reale (POST/PUT)} ([yshift=-2mm]srv.west);
        \draw[arrow, gray, dashed] ([yshift=-4mm]srv.west) -- node[below, black] {4. Risposta (Dati JSON)} ([yshift=-4mm]br.east);
    \end{tikzpicture}
\end{frame}

\begin{frame}{I Cookie e la Sicurezza}
    Memorizzazione dati lato client inviata automaticamente.
    \begin{itemize}
        \item \textbf{HttpOnly}: Protegge da attacchi XSS (JavaScript non può leggerlo).
        \item \textbf{Secure}: Inviato solo su HTTPS.
        \item \textbf{SameSite}: Protegge da attacchi CSRF (Lax, Strict, None).
    \end{itemize}
    \vspace{0.3cm}
    \textit{Importante: Next.js API Routes eliminano spesso i problemi di CORS perché frontend e backend condividono la stessa origine!}
\end{frame}

\section{Integrazione}
\begin{frame}[fragile]{Connessione al Backend}
    Usiamo \texttt{useEffect} per il data fetching all'avvio.
    \begin{lstlisting}[language=javascript]
useEffect(() => {
  fetch("http://localhost:3000/tasks")
    .then(res => res.json())
    .then(data => setTasks(data));
}, []);
\end{lstlisting}
    \textit{Nota: Assicurati che il backend della Lezione 4 supporti il CORS!}
\end{frame}

\begin{frame}{Visione Full-Stack: API Routes}
    Next.js può ospitare anche la logica backend.
    \begin{itemize}
        \item Cartella \texttt{app/api/}: Definisce endpoint HTTP.
        \item Un'unica applicazione per Frontend e Backend.
        \item Niente più problemi di CORS!
    \end{itemize}
    \vspace{0.3cm}
    \centering
    \begin{tikzpicture}
        \node (app) [draw, thick, rounded corners, minimum width=4cm, minimum height=1.5cm, fill=primary!20] {\textbf{Unified Next.js App}};
        \node (fe) [draw, below=0.5cm of app, xshift=-1cm, fill=white] {Frontend};
        \node (be) [draw, below=0.5cm of app, xshift=1cm, fill=white] {API Routes};
        \draw [dashed] (app.south) -- (fe.north);
        \draw [dashed] (app.south) -- (be.north);
    \end{tikzpicture}
\end{frame}

\end{document}
